\section{Numerical values and scope}\label{sec:discusion}

\subsection{Numerical values}

Table~\ref{tab:numerica} reports a deterministic numerical solution of the rescaled Euler system. We use Newton's method on spatial
grids of $64$, $128$, and $256$ points with tolerance $3\cdot10^{-12}$, a grid in $x$ of spacing $0.01$ with additional nodes, natural
spline interpolation, and Simpson's rule for $\int\log F$. The derivative $F'(0)$ is obtained by implicit differentiation of the
discretized system. The maximum difference between the curves for the grids of $256$ and $128$ points, measured at $1001$ nodes, is
$1.085\cdot10^{-5}$. The maximum Newton residual is $2.607\cdot10^{-12}$. The slope obtained by implicit differentiation and the
centered difference with step $10^{-4}$ at $x=0$ on the grid of $256$ points differ by $4.63\cdot10^{-9}$. Nobel, Agrawal, and
Boyd~\cite{NAB} also use Newton's method for the discretized problem at mass one.

\begin{table}[ht] \centering \begin{tabular}{lrrr}
\toprule
Grid & $F(0)$ & $F'(0)$ & $\int_0^1\log F$\\
\midrule
$64$ & $0.08086911055$ & $0.5157082194$ & $-0.9441617254$\\
$128$ & $0.08090292111$ & $0.5157547494$ & $-0.9440411866$\\
$256$ & $0.08091137290$ & $0.5157663849$ & $-0.9440110586$\\
\midrule
Richardson & $0.080914$ & $0.515770$ & $-0.944001$\\
\bottomrule
\end{tabular} \caption{Numerical values of the function $F$. The differences between successive grids decrease by a factor close to
four, consistent with second-order convergence. The last row gives the Richardson extrapolation of that order from the grids of $128$
and $256$ points, calculated from the preceding rows. These are numerical values, not certified by interval arithmetic.} \label{tab:numerica}
\end{table}

By Proposition~\ref{prop:intlogF}, $\sopt(1)=\int_0^1\log F-1\approx-1.9440$. The value $-\sopt(1)\approx1.9440$ lies in the interval
$\alpha=1.942\pm3\cdot10^{-3}$ from~\cite{Simkin}, and the extrapolated value $-\sopt(1)\approx1.944001$ lies in the interval
$[1.944000752,\,1.944001082]$ from~\cite{NAB}. Remark~\ref{obs:simkin} explains the normalization that allows the two quantities to be
compared. The agreement of the decimal digits serves as a numerical check. The extrapolations using grids $(64,128)$ and $(128,256)$
differ by less than $10^{-8}$. The quadrature error in $x$ has not been estimated. \ifx\DOIcodigo\empty\else A program that reproduces
these values by another method, a fixed-point iteration, is available in~\cite{Codigo}. Its results differ from the table by less than
$10^{-8}$. \fi Using Monte Carlo simulation on boards of size $1024\times1024$, Liu, Liao, and Wang~\cite{LLW} estimate
$\alpha\approx1.946\pm0.003$ by integrating the specific heat starting at infinite temperature.

\subsection{Interpretation at the full board}

When $x=0$, the row and column vacancies vanish. The factors $r_{R,0}$ and $r_{C,0}$ describe the limiting row and column vacancy
densities, respectively, each normalized to mass one. Together with the limiting diagonal vacancies and the factor $b_0$, they
determine $\nu_1$ by Theorem~\ref{thm:lleno}. The spatial regularity obtained is Lipschitz, with a common bound for $0\le x\le 1/2$.
Analyticity describes a different dependence, that of the profiles on the parameter $x$.

All statements in this paper concern the continuous problem. The constants in Sections~\ref{sec:cotas} and~\ref{sec:regularidad} are
explicit but too conservative for numerical use. Table~\ref{tab:numerica} estimates $F(0)$ by solving the continuous system. The
identity $F(0)=1/b_0$ follows from Theorem~\ref{thm:F}. Interpreting this value as a frequency on finite boards also requires
specifying the probability law on placements and justifying the passage from the discrete model to the continuous one.
